\documentclass[11pt]{article}
\usepackage{amsmath,amssymb,amsthm,mathtools}
\usepackage[margin=1.1in]{geometry}
\usepackage{booktabs}
\theoremstyle{plain}
\newtheorem{theorem}{Theorem}
\newtheorem{proposition}[theorem]{Proposition}
\newtheorem{corollary}[theorem]{Corollary}
\newtheorem{lemma}[theorem]{Lemma}
\theoremstyle{definition}
\newtheorem{hypothesis}[theorem]{Hypothesis}
\newtheorem{remark}[theorem]{Remark}

\title{The vertical migration step is a unit translation,\\
and the orientation gap is irreducible}
\author{Clio}
\date{6 October 2026}

\begin{document}
\maketitle

\begin{abstract}
In \texttt{proofs/2026-10-05-c3-migration-length-grading.tex} the migration
length $\ell(M)$ of Purbhoo's mosaic migration (\texttt{0705.1184}) was shown to
admit an exact potential $\varphi(a,b,c,e)=b+e$ on the $\mathbb{Z}^4$ lift, with
$\Delta\varphi=1$ at every step --- except that the strictly vertical
$180^\circ$ step exists in two directions, and $91$ of $707$ computed steps were
strictly vertical with all $91$ taking $\Delta\varphi=+1$ and no proof excluding
$-1$.  This note sharpens that gap.  We show the two vertical steps are one
geometric picture read in two directions (Proposition~\ref{prop:slide}), deduce
that \emph{no} choice of potential can assign $+1$ to both
(Corollary~\ref{cor:irred}), so the gap can only be closed by a reachability
argument; we show hexagon minimality cannot discriminate
(Proposition~\ref{prop:minimality}); and we widen and, more importantly,
\emph{explain} the empirical null, from $707$ to $3111$ steps.  The gap is not
closed.  Two hypotheses on which Theorem~1 of the earlier paper rests are
re-examined, and the status of the second is corrected.
\end{abstract}

\section{Setting}

Vertices of a mosaic are integer $4$-tuples $(a,b,c,e)$ standing for
$a\,u_0+b\,u_{60}+c\,u_{30}+e\,u_{90}$, where $u_t$ is the unit vector at angle
$t$; the map $\mathbb{Z}^4\to\mathbb{R}^2$ is injective.  Set
$\varphi(a,b,c,e)=b+e$, extended to a tile as the mean of its four vertices.
Migration travels in the direction $u=u_{150}$ (the north direction $N_A$ of nest
$A$), and we write $\sigma=\langle\,\cdot\,,u_{150}\rangle$ for the travel
coordinate.  Purbhoo's step rotates the smallest hexagon containing the
travelling rhombus and lying weakly right of its leftmost point, by $180^\circ$
(centrally symmetric hexagon) or $\pm120^\circ$ (three-fold symmetric hexagon).
A step is \emph{strictly vertical} when $\sigma$ is unchanged.

\section{The vertical step is a unit translation}

\begin{proposition}\label{prop:slide}
Up to translation there are exactly two strictly vertical migration steps, and
they share the same hexagon region
\[
   H \;=\; R \oplus [0,1]\,u_{60},
\]
the Minkowski sum of the travelling rhombus $R$ with a unit segment
perpendicular to the travel direction.  $H$ is the zonogon on
$\{u_0,u_{60},u_{150}\}$, has four tiles, and has exactly two tilings,
interchanged by its central symmetry.  The $180^\circ$ rotation translates the
rhombus by $+u_{60}$ in one tiling and by $-u_{60}$ in the other, giving
$\Delta\varphi=\varphi(u_{60})=+1$ and $-1$ respectively.
\end{proposition}

\begin{proof}
Exhaustive enumeration of every admissible hexagon up to translation ($90$
regions), every tiling, every rhombus and every admissible rotation, filtered by
Purbhoo's constraints (the rhombus attains the hexagon's $\sigma$-minimum; the
rotated rhombus keeps one of the two orientations of Proposition 3.1), returns
exactly two steps with $\sigma$ unchanged.  Both have the hexagon with vertex
cycle $(0,0,-1,1),(0,1,-1,1),(1,1,-1,1),(1,1,0,0),(1,0,0,0),(0,0,0,0)$, i.e.\ the
zonogon on $u_0,u_{60},u_{150}$: of its three pairs of edge directions,
$\{u_0,u_{60}\}$ spans a $60^\circ$ rhombus (two triangles), $\{u_0,u_{150}\}$
the $30/150$ rhombus tile and $\{u_{60},u_{150}\}$ a square, so $|H|=4$ tiles.
In one tiling the rhombus is $R=\{(0,0,-1,1),(0,0,0,0),(1,0,0,0),(1,0,-1,1)\}$
and in the other it is $R+(0,1,0,0)=R+u_{60}$, and these are the two tilings of a
three-direction zonogon.  Since $u_{60}\perp u_{150}$, both have the same
$\sigma$-extent, which is why the step is vertical; and
$\varphi$ is $\mathbb{Z}$-linear with $\varphi(u_{60})=1$, so
$\Delta\varphi=\pm1$ according to the direction.
\end{proof}

\begin{corollary}\label{cor:irred}
There is no function $\psi$ on rhombus positions with $\Delta\psi=+1$ on both
vertical steps.  Hence the gap cannot be closed by replacing $\varphi$ with a
better potential; it can be closed only by showing that one of the two
directions is never realised.
\end{corollary}

\begin{proof}
By Proposition~\ref{prop:slide} the two steps are mutually inverse: the same
region, the same rotation, the two tilings.  So for any $\psi$,
$\Delta\psi(T_+\!\to\!T_-)=-\Delta\psi(T_-\!\to\!T_+)$, and the two cannot both
equal $+1$.
\end{proof}

\begin{proposition}\label{prop:minimality}
Minimality of the hexagon does not discriminate between the two directions.
\end{proposition}

\begin{proof}
Enumerating every admissible local move for the canonical nest-$A$ rhombus gives
$15$ moves, and \emph{every one of them} uses a four-tile hexagon:
one vertical with $\Delta\varphi=-1$, one vertical with $+1$, one forward
two-fold with $+1$, six backward three-fold with $-1$ (excluded by forward
motion) and six forward three-fold with $+1$.  Since all candidate hexagons have
the same size, ``the smallest hexagon containing the rhombus'' cannot prefer one
vertical direction over the other.  Note also that $H=R\oplus[0,1]u_{60}$ has the
same $\sigma$-extent as $R$, so the clause ``lying weakly right of the rhombus's
leftmost point'' is satisfied by \emph{both} strips and is likewise silent.
\end{proof}

So the asymmetry, if any, is global.  Proposition~\ref{prop:slide} does localise
it to a concrete statement: the $180^\circ$ rotation maps the strip's rear end to
its front end, so the two strips differ in where the square sits relative to
travel.

\begin{hypothesis}[G]\label{hyp:G}
In a mosaic the travelling rhombus never has a free unit strip on its
$-u_{60}$ side; equivalently, it never has the configuration
\emph{(square ahead, two triangles)} occupying $R\ominus[0,1]u_{60}$.
\end{hypothesis}

\section{The null, widened and explained}

The earlier figure was $91$ of $707$ steps strictly vertical, all $+1$.  That
count alone was not evidence.  The step engine breaks a tie among equal-progress
candidates by maximising progress, and progress is \emph{constant} on vertical
candidates, so the tie-break returns whichever candidate the enumeration happened
to list first.  Had the $-1$ slide ever been in the candidate pool, an all-$+1$
tally would have been an artifact of enumeration order and not a fact about
mosaics.  So availability was measured, not selection, over
$(n,d)=(4,2),(5,2),(5,3),(6,3)$: $362$ mosaics, $3111$ steps.

\begin{center}
\begin{tabular}{lr}
\toprule
steps logged & 3111\\
chosen step strictly vertical & 434\\
\quad of those, $\Delta\varphi=+1$ & 434\\
\quad of those, $\Delta\varphi=-1$ & 0\\
\midrule
steps where a vertical $\Delta\varphi=-1$ candidate was \textbf{available} & \textbf{0}\\
steps where a vertical $\Delta\varphi=+1$ candidate was available & 434\\
steps where both were available & 0\\
\bottomrule
\end{tabular}
\end{center}

The silence is \emph{absence}, not deselection: the tie-break is never consulted
on a vertical step, so the null is a statement about mosaic geometry rather than
about the implementation.  Hypothesis~\ref{hyp:G} is therefore an empirical null
at $3111$ steps and $362$ mosaics, and remains unproved.

\section{A second hypothesis, and a correction to what I claimed about it}

\emph{Correction first.}  I set out believing this hypothesis was unnamed, hidden
in a filter of the classification script.  It is not.  The earlier paper states it
as condition (C2) of its classification lemma --- ``the hexagon contains exactly one
rhombus'' --- with the justification ``during a migration all other rhombi are
packed in nests'', and it already records the configurations this condition
excludes (``without (C1)--(C3) the classification admits $210$ forward
configurations with $\Delta\phi\in\{0,\pm\tfrac12\}$'').  So the $\pm\tfrac12$
phenomenon below is \emph{not} a discovery of this session either.  What is new is
narrower, and is the measurement and the status, not the statement.

(C2)'s justification is an assertion, not a proof; and the step engine imposes no
such restriction, so (C2) is a property of the \emph{chosen} hexagon rather than a
consequence of the stated rule.  Measuring what the engine actually does:

\begin{center}
\begin{tabular}{lr}
\toprule
chosen steps & 3111\\
with $\Delta\varphi=1$ (bare height, no type correction) & 3111\\
whose chosen hexagon contains exactly one rhombus & 3111\\
whose chosen hexagon has exactly four tiles & 3111\\
\bottomrule
\end{tabular}
\end{center}

So the classification does cover every chosen step in this range --- but
``exactly one rhombus'' is a property of the \emph{chosen} hexagon, not a
consequence of the stated rule, and the potential identity needs it:

\begin{hypothesis}[G2 $=$ (C2)]\label{hyp:G2}
The minimal admissible hexagon of the travelling rhombus contains no other
rhombus.
\end{hypothesis}

So the classification does cover every chosen step here, and (C2) holds on all of
them --- but as a measured fact, not a proved one.  It is load-bearing, because the
candidate pool \emph{did} contain moves violating the potential identity: $173$
candidates with bare $\Delta\varphi=+\tfrac12$ and $173$ with $-\tfrac12$.  They
were always rejected, and the reason is now established rather than asserted: all $346$ rejected
candidates had rhombus orientation $(120,90)$, which Purbhoo's Proposition 3.1
forbids, and all $346$ had the horizontal/vertical flag false, while all $346$
chosen moves had it true.  Both controls returned zero (no rejected candidate had
an allowed orientation; no chosen candidate had a forbidden one).  Hence:

\begin{proposition}
On these $3111$ steps, Purbhoo's horizontal/vertical tie-break selects exactly
the orientations permitted by his Proposition 3.1.  The tie-break is not an
independent heuristic; it implements the orientation restriction.
\end{proposition}

A byproduct: no rhombus-type correction is needed in the potential.  The bare
height $\varphi$ has $\Delta\varphi=1$ on every chosen step, all $3111$.

Finally, the move repertoire of a journey: $1130$ two-fold steps
$(0,150)\to(0,150)$, $696$ two-fold $(60,30)\to(60,30)$, $1112$ three-fold
$(0,150)\to(60,30)$ and $173$ three-fold $(60,30)\to(0,150)$.  A journey slides
within an orientation by $180^\circ$ rotations and alternates orientation by
three-fold rotations, starting in nest $A$'s orientation (spanned by $E_A=u_0$,
$N_A=u_{150}$) and ending in nest $B$'s (spanned by $E_B=u_{240}$, $N_B=u_{30}$).

\section*{Gaps}

\begin{enumerate}
\item Hypotheses~\ref{hyp:G} and \ref{hyp:G2} are unproved.  Theorem~1 of
  \texttt{2026-10-05-c3-migration-length-grading.tex} must state both as
  hypotheses.  (G2) is already there as condition (C2) of its classification
  lemma, justified by an assertion; (G) is there as its stated Gap.  What this note
  changes is that (G) is now known to be \emph{irreducible}
  (Corollary~\ref{cor:irred}) rather than merely open, and that the evidence for it
  is availability rather than selection.
\item I did not amend the earlier paper.  Its Theorem~1 and Gap should be rewritten
  to cite Corollary~\ref{cor:irred} and the $3111$-step availability measurement,
  and to say that (C2) is a hypothesis about the chosen hexagon rather than a
  consequence of the rule.  That edit is outstanding.
\item Corollary~\ref{cor:irred} rules out the entire family of repairs by a
  better potential, so no further search in that direction is warranted.
\item The reach is unchanged on one axis: no fibre with
  $c^\lambda_{\mu\nu}\ge3$ was run (the tiler is exponential);
  $c\ge2$ was exercised on $21$ fibres.  The $3111$-step sweep is not exhaustive
  in the multiplicity direction.
\end{enumerate}

\section*{Apparatus}

\texttt{work-in-progress/code/mosaic-migration-1005c3/} (unchanged), driven by
the measurement scripts of this session.

\end{document}
